% Options for packages loaded elsewhere
\PassOptionsToPackage{unicode}{hyperref}
\PassOptionsToPackage{hyphens}{url}
\documentclass[
  11pt,
  letterpaper,
]{article}
\usepackage{xcolor}
\usepackage[margin=1in]{geometry}
\usepackage{amsmath,amssymb}
\setcounter{secnumdepth}{-\maxdimen} % remove section numbering
\usepackage{iftex}
\ifPDFTeX
  \usepackage[T1]{fontenc}
  \usepackage[utf8]{inputenc}
  \usepackage{textcomp} % provide euro and other symbols
\else % if luatex or xetex
  \usepackage{unicode-math} % this also loads fontspec
  \defaultfontfeatures{Scale=MatchLowercase}
  \defaultfontfeatures[\rmfamily]{Ligatures=TeX,Scale=1}
\fi
\usepackage{lmodern}
\ifPDFTeX\else
  % xetex/luatex font selection
\fi
% Use upquote if available, for straight quotes in verbatim environments
\IfFileExists{upquote.sty}{\usepackage{upquote}}{}
\IfFileExists{microtype.sty}{% use microtype if available
  \usepackage[]{microtype}
  \UseMicrotypeSet[protrusion]{basicmath} % disable protrusion for tt fonts
}{}
\makeatletter
\@ifundefined{KOMAClassName}{% if non-KOMA class
  \IfFileExists{parskip.sty}{%
    \usepackage{parskip}
  }{% else
    \setlength{\parindent}{0pt}
    \setlength{\parskip}{6pt plus 2pt minus 1pt}}
}{% if KOMA class
  \KOMAoptions{parskip=half}}
\makeatother
\usepackage{graphicx}
\makeatletter
\newsavebox\pandoc@box
\newcommand*\pandocbounded[1]{% scales image to fit in text height/width
  \sbox\pandoc@box{#1}%
  \Gscale@div\@tempa{\textheight}{\dimexpr\ht\pandoc@box+\dp\pandoc@box\relax}%
  \Gscale@div\@tempb{\linewidth}{\wd\pandoc@box}%
  \ifdim\@tempb\p@<\@tempa\p@\let\@tempa\@tempb\fi% select the smaller of both
  \ifdim\@tempa\p@<\p@\scalebox{\@tempa}{\usebox\pandoc@box}%
  \else\usebox{\pandoc@box}%
  \fi%
}
% Set default figure placement to htbp
\def\fps@figure{htbp}
\makeatother
\setlength{\emergencystretch}{3em} % prevent overfull lines
\providecommand{\tightlist}{%
  \setlength{\itemsep}{0pt}\setlength{\parskip}{0pt}}
\usepackage{bookmark}
\IfFileExists{xurl.sty}{\usepackage{xurl}}{} % add URL line breaks if available
\urlstyle{same}
\hypersetup{
  hidelinks,
  pdfcreator={LaTeX via pandoc}}

\author{}
\date{}

\begin{document}

\section{Elinor Ostrom's Design Principles and the Tholonic N-D-C
Framework: Governance as Structural
Balance}\label{elinor-ostroms-design-principles-and-the-tholonic-n-d-c-framework-governance-as-structural-balance}

\textbf{Author:} J. W. Milton, Clarity Coalition

\textbf{Version:} 1.0

\textbf{Date:} 26 June 2026

\textbf{Keywords:} Elinor Ostrom; common-pool resources; design
principles; tholonic model; N-D-C framework; tragedy of the commons;
polycentric governance; institutional analysis; sustainability;
governance balance score

\begin{center}\rule{0.5\linewidth}{0.5pt}\end{center}

\subsection{Abstract}\label{abstract}

Elinor Ostrom's Nobel Prize-winning work on the governance of
common-pool resources (CPR) identified eight design principles that
reliably distinguish self-governing institutions that sustain shared
resources from those that collapse. These principles were derived
empirically from field data across dozens of cultures and resource
types: alpine meadows in Switzerland, fisheries in Turkey and Japan,
irrigation systems in Spain and the Philippines. This paper argues that
Ostrom's eight principles are not an arbitrary empirical checklist but a
structural finding: they describe the conditions under which a CPR
governance institution achieves and maintains the N-D-C balance
predicted by the tholonic model.

The paper makes four claims. First, the eight design principles divide
cleanly into D-type principles (boundary definition, rule specification,
monitoring, graduated sanctions, conflict resolution), which constitute
the constraining apparatus of the governance tholon, and C-type
principles (collective-choice arrangements, external recognition of
rights, nested enterprises), which constitute its integrating and
participatory apparatus. Second, Ostrom's documented failure mode, the
tragedy of the commons, is a C-dominant partial tholon: collective
extraction (C) is active while institutional constraint (D) is absent or
too weak to match it. Third, a structurally distinct second failure
mode, high-D regulatory collapse, occurs when D dominates without C:
centralized top-down regulation without collective participation
produces rigid, brittle institutions that cannot adapt. Ostrom's
framework identifies both failure types but does not have a formal
vocabulary for the distinction. The tholonic model provides it. Fourth,
the tholonic balance score \(B\) provides a continuous governance health
metric that extends Ostrom's binary principle-presence assessment to a
graded scalar, enabling differentiated predictions about partial
institutions.

Illustrative B-scores are computed for four cases drawn from Ostrom's
documented case studies: Törbel, Switzerland (near-complete
institutional balance, \(B \approx 100\)); Alanya, Turkey
(above-threshold partial balance, \(B \approx 67\)); open-access fishery
collapse (C-dominant failure, \(B \to 0\)); and Atlantic groundfish
regulatory failure (D-dominant failure, \(B \to 0\) via the opposite
imbalance). Ostrom's framework predicts that all three non-ideal cases
are problematic; the tholonic model additionally predicts that the
C-dominant and D-dominant failures will produce structurally different
collapse patterns with different time profiles and different
intervention requirements. Three falsifiable predictions are stated.

\begin{center}\rule{0.5\linewidth}{0.5pt}\end{center}

\subsection{1. Introduction}\label{introduction}

\textbf{What this paper provides.} A systematic mapping of Ostrom's
eight design principles onto tholonic D and C roles; a formal account of
why the tragedy of the commons is a C-dominant partial tholon; the
identification of a structurally distinct D-dominant failure mode not
captured by Ostrom's framework; illustrative B-scores for four
documented case studies; three falsifiable predictions that distinguish
the tholonic account from Ostrom's framework; and a connection to the
TVPCI supply chain transparency scoring model developed in earlier
papers of this series.

\textbf{What this paper does not provide.} Original field data on CPR
governance; a comprehensive survey of Ostrom's complete body of work; a
claim that Ostrom was unaware of the underlying structural logic (she
may well have intuited elements of it); or primary statistical analysis
of Ostrom's case study database. The B-scores computed here are
illustrative, derived from the qualitative principle-presence
assessments in Ostrom's own published case studies, not from independent
measurement.

\textbf{Organization.} §2 introduces Ostrom's framework and the CPR
problem. §3 briefly recaps the tholonic N-D-C model. §4 maps the CPR
governance system onto the tholonic framework and classifies the eight
principles. §5 analyzes four case studies with B-scores. §6 identifies
what the tholonic model adds beyond Ostrom's framework. §7 states
falsifiable predictions. §8 connects the governance analysis to the
project's supply chain work. §9 concludes.

\begin{center}\rule{0.5\linewidth}{0.5pt}\end{center}

\subsection{2. Ostrom's Framework: Common-Pool Resources and the Design
Principles}\label{ostroms-framework-common-pool-resources-and-the-design-principles}

\subsubsection{2.1 The Common-Pool Resource
Problem}\label{the-common-pool-resource-problem}

A common-pool resource (CPR) is a resource system characterized by two
properties {[}Ost90{]}:

\textbf{Subtractability}: each unit consumed by one user reduces the
stock available to others. A fish caught is not available to the next
fisher; water extracted from an aquifer lowers the water table for all
users; a pasture grazed reduces the forage available next season.

\textbf{Difficulty of exclusion}: it is costly or impossible to prevent
potential users from accessing the resource. Open fisheries, shared
aquifers, common pastures, and unenclosed forests are all CPRs: barriers
to entry are high to enforce or absent entirely.

This combination produces the structural problem that Garrett Hardin
{[}Har68{]} formalized as the ``tragedy of the commons'': if each
individual user acts rationally to maximize personal extraction, the
cumulative effect depletes the resource to the detriment of all. The
dominant theoretical prescriptions were either privatize the resource
(align individual incentives with conservation through property rights)
or regulate it through central state authority.

Ostrom's contribution was to show, through decades of field research
spanning more than fifty CPR case studies across six continents, that a
third path existed and had been practiced successfully for centuries in
many communities without awareness of either prescription {[}Ost90{]}.
Self-governing communities regularly achieve sustainable resource
management through locally evolved institutional rules, and the
institutions that succeed share a recognizable structural pattern.

\subsubsection{2.2 The Eight Design
Principles}\label{the-eight-design-principles}

From the comparative analysis of successful and failed CPR governance
institutions, Ostrom identified eight design principles (she later
refined these with colleagues, but the eight-principle formulation from
\emph{Governing the Commons} {[}Ost90{]} is the canonical statement):

\textbf{Principle 1: Clearly defined boundaries.} Individuals or
households who have rights to withdraw resource units from the CPR must
be clearly defined, as must the boundaries of the CPR itself.

\textbf{Principle 2: Congruence between rules and local conditions.}
Appropriation and provision rules are matched to local conditions,
including the relevant attributes of the resource system and the
community.

\textbf{Principle 3: Collective-choice arrangements.} Most individuals
affected by the operational rules can participate in modifying those
rules.

\textbf{Principle 4: Monitoring.} Monitors who actively audit CPR
conditions and appropriator behavior are present, and are accountable to
the appropriators or are themselves appropriators.

\textbf{Principle 5: Graduated sanctions.} Appropriators who violate
operational rules are likely to receive graduated sanctions (ranging
from very low to very high, depending on the seriousness and context of
the offense) from other appropriators, from officials accountable to
these appropriators, or from both.

\textbf{Principle 6: Conflict-resolution mechanisms.} Appropriators and
their officials have rapid access to low-cost local arenas to resolve
conflicts among appropriators or between appropriators and officials.

\textbf{Principle 7: Minimal recognition of rights to organize.} The
rights of appropriators to devise their own institutions are not
challenged by external governmental authorities.

\textbf{Principle 8: Nested enterprises} (for systems embedded in larger
systems). Appropriation, provision, monitoring, enforcement, conflict
resolution, and governance activities are organized in multiple layers
of nested enterprises.

\subsubsection{2.3 The Evidential Weight of These
Principles}\label{the-evidential-weight-of-these-principles}

A crucial point for the tholonic analysis is that these principles were
not derived from theory and then tested against data. They were
extracted from the data first, through systematic comparison of case
studies, and the theoretical interpretation followed. Their empirical
provenance means that any structural model that can account for why
these eight and not some other set should be considered structurally
significant, not merely descriptive.

\subsubsection{2.4 Polycentric Governance}\label{polycentric-governance}

Beyond the eight principles, Ostrom developed the concept of polycentric
governance {[}Ost10{]}: the observation that large-scale CPR systems are
most reliably managed by multiple overlapping centers of decision-making
authority at different scales, each with jurisdiction appropriate to the
scale of the resource dynamics it governs. This is not a normative
prescription but an empirical regularity: monocentric governance,
whether private or governmental, consistently underperforms polycentric
governance for CPRs at regional and global scales.

\begin{center}\rule{0.5\linewidth}{0.5pt}\end{center}

\subsection{3. The Tholonic N-D-C
Framework}\label{the-tholonic-n-d-c-framework}

The tholonic model is developed in full in papers
\href{https://github.com/baardev/truevalue/raw/main/docnav/Research/papers/1_recursive-tholonic-five-constants/1_recursive-tholonic-five-constants.pdf}{1}
through
\href{https://github.com/baardev/truevalue/raw/main/docnav/Research/papers/3_minimal-recursive-triadic-framework/3_minimal-recursive-triadic-framework.pdf}{3}
of this series
\href{https://github.com/baardev/truevalue/raw/main/docnav/Research/papers/1_recursive-tholonic-five-constants/1_recursive-tholonic-five-constants.pdf}{Mil26a},
\href{https://github.com/baardev/truevalue/raw/main/docnav/Research/papers/2_supply-chain-transparency-tvpci/2_supply-chain-transparency-tvpci.pdf}{Mil26b},
\href{https://github.com/baardev/truevalue/raw/main/docnav/Research/papers/3_minimal-recursive-triadic-framework/3_minimal-recursive-triadic-framework.pdf}{Mil26c}.
The minimal recap required here is as follows.

A \emph{tholon} is any stable, self-sustaining structure simultaneously
instantiating three functional roles. \textbf{N (Negotiation)} is the
emergent stable equilibrium: the coherent state that persists as a
consequence of D and C remaining in approximate balance. \textbf{D
(Definition)} is the constraining apparatus: all mechanisms that limit,
bound, specify, and filter. D is internally focused. \textbf{C
(Contribution)} is the integrating-expressive apparatus: all mechanisms
that generate, connect, and produce output. C is externally focused.

A \emph{partial tholon} realizes only one or two of these roles and is
structurally unstable. D-dominant systems become rigid and brittle.
C-dominant systems dissipate without coherence. The balance score is:

\[B(D,C) = \frac{2 \cdot \min(D,C)}{D + C} \cdot 100\]

The stability threshold is \(B = 61.8\) (\(1/\varphi\), where
\(\varphi\) is the golden ratio). Below this threshold, the N-state is
structurally fragile. The relation between N, D, and C is recursive: a
parent N differentiates into D and C, which negotiate to produce a child
N, which becomes the parent N of the next level.

\begin{center}\rule{0.5\linewidth}{0.5pt}\end{center}

\subsection{4. Mapping Ostrom onto the Tholonic
Framework}\label{mapping-ostrom-onto-the-tholonic-framework}

\subsubsection{4.1 The CPR System as a
Tholon}\label{the-cpr-system-as-a-tholon}

The CPR governance system maps onto the tholonic framework at two levels
simultaneously.

\textbf{Level 1: The resource.} The resource system itself (the fishery,
the aquifer, the alpine meadow) is the N-state at the ecological level.
It is the stable, self-sustaining state that the governance institution
exists to maintain. The resource's D component is its natural
regeneration rate and ecological constraints; its C component is the
harvest and use flows that human communities extract from it. A
sustainable resource is one where \(D \approx C\): extraction
approximately matches regeneration. Ostrom calls this the ``sustainable
yield'' condition; the tholonic model formalizes it as ecological
N-balance.

\textbf{Level 2: The governance institution.} The institution itself is
a tholon whose N-state is the stable governance configuration (the set
of rules that is actually followed, respected, and enforced over time).
Its D component is all the institutional mechanisms that constrain and
define extraction (boundary rules, provision rules, monitoring,
sanctions, conflict resolution). Its C component is all the
institutional mechanisms that integrate participation and connect the
institution to its environment (collective choice, external recognition,
nested enterprises).

\begin{figure}
\centering
\pandocbounded{\includegraphics[keepaspectratio]{/home/jw/src/tv/docnav/Research/papers/17_ostrom-tholonic-governance/figures/17_cpr-tholon-triangle.png}}
\caption{N-D-C triangle diagram showing the CPR governance tholon: N at
top as the stable governance configuration and sustainable resource, D
lower-right as institutional constraints and boundary rules, C
lower-left as collective participation and integrative mechanisms.}
\end{figure}

The two levels are coupled. The governance institution (Level 2) exists
to maintain the resource (Level 1). If the governance institution
achieves N-balance (the rules are followed, respected, and adaptive),
the resource level also tends toward N-balance (sustainable yield). If
the governance institution collapses into a partial tholon, the resource
collapses in a predictable direction depending on which component (D or
C) is missing.

\subsubsection{4.2 D-Type and C-Type
Principles}\label{d-type-and-c-type-principles}

The eight design principles divide cleanly by tholonic role:

\textbf{D-type principles (constraining apparatus):}

\begin{itemize}
\tightlist
\item
  \emph{Principle 1 (Defined boundaries):} definitional by construction.
  It specifies who is and is not inside the resource system and the
  institution. Establishing identity and limits is the primary D
  function.
\item
  \emph{Principle 2 (Rules match local conditions):} the specification
  of what rules apply under what local constraints. Rule specificity is
  a D function.
\item
  \emph{Principle 4 (Monitoring):} the measurement and enforcement of
  compliance with rules. Monitoring is the operational expression of the
  constraining apparatus: it makes rules real rather than nominal.
\item
  \emph{Principle 5 (Graduated sanctions):} the enforcement mechanism
  that gives consequences to rule violation. Sanctions are D's teeth.
\item
  \emph{Principle 6 (Conflict resolution):} mechanisms that resolve
  disputes about rule interpretation and application. Conflict
  resolution maintains D-coherence: it ensures that the constraining
  apparatus does not fragment into competing claims.
\end{itemize}

\textbf{C-type principles (integrating-participatory apparatus):}

\begin{itemize}
\tightlist
\item
  \emph{Principle 3 (Collective-choice arrangements):} the mechanism by
  which those affected by rules participate in modifying them.
  Collective choice is integrative: it connects the lived experience of
  resource users (C-signals from the resource level) to the rule
  structure (D). Without this C principle, D becomes rigid.
\item
  \emph{Principle 7 (External recognition of rights):} the connection of
  the institution to the larger governance environment. External
  recognition is the outward-facing, integrative function that embeds
  the local institution within the systems around it. Without this, even
  a well-balanced local institution is fragile to external disruption.
\item
  \emph{Principle 8 (Nested enterprises):} the hierarchical structure of
  governance at multiple scales. Nesting is the explicit C-D scaling
  mechanism: it ensures that C flows at the local level connect to D
  structures at the regional level, and so on. This is the tholonic
  recursive hierarchy stated as a governance principle.
\end{itemize}

\begin{figure}
\centering
\pandocbounded{\includegraphics[keepaspectratio]{/home/jw/src/tv/docnav/Research/papers/17_ostrom-tholonic-governance/figures/17_principles-classification.png}}
\caption{Grouped bar chart showing the eight design principles
classified by D or C role, with illustrative presence scores for four
case studies: Törbel Switzerland, Alanya Turkey, open-access failure,
and D-dominant regulatory failure.}
\end{figure}

\subsubsection{4.3 The Tragedy of the Commons as C-Dominant Partial
Tholon}\label{the-tragedy-of-the-commons-as-c-dominant-partial-tholon}

Hardin's tragedy scenario {[}Har68{]} describes an institution with zero
D and non-zero C. Every individual exercises their capacity to extract
(C-behavior: contribution of individual extraction to the collective
pool, but from the institution's perspective, pure C-runaway without
constraint). No rules define who may extract, in what quantity, with
what monitoring or sanction. The institution has no D apparatus at all.

In tholonic terms this is a pure C-dominant partial tholon:
\(D \approx 0\), \(C > 0\). The balance score is \(B \to 0\). The
N-state (sustainable resource management) is impossible because N cannot
form without D. Resource collapse is not a failure of individual
rationality; it is the structural consequence of attempting to maintain
N with only one of its two constitutive components.

Ostrom's empirical correction to Hardin, that communities often develop
D structures spontaneously to prevent this collapse, is the tholonic
claim that systems under pressure to maintain N tend toward D-C balance.
The community converges on the governance institution not because any
individual chooses the institution over extraction, but because the
structural attractor of a system under N-pressure is balance, not either
extreme.

\subsubsection{4.4 The D-Dominant Failure Mode: Regulatory
Brittleness}\label{the-d-dominant-failure-mode-regulatory-brittleness}

Ostrom's framework identifies the tragedy of the commons (C-dominant) as
the primary failure mode. It does not provide an equally formal account
of the opposite failure: D-dominant governance collapse. The tholonic
model predicts this failure mode as clearly as the C-dominant one, and
it has a different structural profile.

A D-dominant governance institution is one where boundary rules,
monitoring, and sanctions are strong and centralized, but
collective-choice arrangements and external integration are absent or
suppressed. Top-down regulatory management of fisheries by national
governments exemplifies this: the D apparatus (quotas, licensing,
patrol) is present and formally rigorous, but Principle 3 (collective
choice) and Principle 8 (nested enterprises connecting local knowledge
to regulatory design) are absent.

The tholonic prediction for D-dominant partial tholons is rigidity and
brittleness. The institution cannot adapt to changing conditions because
there is no C mechanism (collective participation) through which new
information flows from resource users to rule-makers. The rules become
outdated relative to the resource state; the institution defends
outdated rules with increasing force; and when the mismatch between
rules and reality becomes critical, the institution collapses suddenly
rather than adapting gradually.

The Atlantic groundfish (cod) collapse of the early 1990s is a
documented case of precisely this pattern. Canadian and U.S. federal
fisheries management had strong D (quotas, monitoring, enforcement) but
suppressed local collective-choice mechanisms (Principle 3) and operated
at a single centralized scale rather than through nested enterprises
(Principle 8). Scientific assessments that underestimated the stock,
combined with politically motivated quota decisions that overrode
monitoring data, produced a D-dominant failure: the institution defended
its quotas past the point at which adaptation was possible, and the
fishery collapsed to near-zero within three years.

\begin{center}\rule{0.5\linewidth}{0.5pt}\end{center}

\subsection{5. Case Studies with B-Score
Analysis}\label{case-studies-with-b-score-analysis}

B-scores are computed by scoring each case study on the presence and
quality of each principle (0 = absent, 0.5 = partially present, 1.0 =
fully present), summing the D-type scores (max 5.0) and C-type scores
(max 3.0), normalizing each to a 0-to-100 scale, and applying the
balance formula. The qualitative assessments are drawn from Ostrom's
documented case studies in \emph{Governing the Commons} {[}Ost90{]} and
subsequent work.

\[D_{\text{norm}} = \frac{\sum \text{D principles scored}}{5} \times 100, \quad C_{\text{norm}} = \frac{\sum \text{C principles scored}}{3} \times 100\]

\[B = \frac{2 \cdot \min(D_{\text{norm}}, C_{\text{norm}})}{D_{\text{norm}} + C_{\text{norm}}} \times 100\]

\subsubsection{5.1 Törbel, Switzerland (Alpine Meadow, Documented Since
1483)}\label{tuxf6rbel-switzerland-alpine-meadow-documented-since-1483}

The Törbel village commons is one of Ostrom's most thoroughly documented
long-run success cases {[}Ost90, Ch. 3{]}. The community has maintained
alpine meadows, forests, and irrigation systems collectively for over
five centuries through formal statutes (the earliest surviving document
dates to 1483) and a village assembly that meets annually.

\textbf{D principles:} Principle 1 (clearly defined boundaries,
documented since 1483): 1.0. Principle 2 (rules match local conditions,
including seasonal grazing calendars calibrated to snowpack): 1.0.
Principle 4 (monitoring by elected guards): 1.0. Principle 5 (graduated
fines for violations, documented in village records): 1.0. Principle 6
(village assembly as conflict resolution forum): 1.0. D-type sum:
5.0/5.0.

\textbf{C principles:} Principle 3 (annual village assembly with all
resident households): 1.0. Principle 7 (Swiss cantonal law recognizes
village commons institutions): 1.0. Principle 8 (nested within cantonal
and federal resource governance): 1.0. C-type sum: 3.0/3.0.

\[D_{\text{norm}} = 100, \quad C_{\text{norm}} = 100, \quad B = 100\]

This perfect B-score reflects an institution that has had five centuries
to achieve full D-C integration. It is structurally coherent in every
dimension the tholonic model identifies.

\subsubsection{5.2 Alanya, Turkey (Inshore Fishery, Self-Organized from
the
1970s)}\label{alanya-turkey-inshore-fishery-self-organized-from-the-1970s}

The Alanya fishery is a case Ostrom presents as a self-organized success
that emerged within a generation {[}Ost90, Ch. 3{]}. Turkish inshore
fishers developed a lottery system for rotating fishing spots,
eliminating competition and overfishing. The institution is genuine and
has functioned for decades.

\textbf{D principles:} Principle 1 (spot boundaries clearly defined,
lottery positions publicly registered): 1.0. Principle 2 (seasonal
calendar matched to fish migration): 1.0. Principle 4 (mutual
monitoring, each fisher observes others at adjacent spots): 1.0.
Principle 5 (social sanction and local fine system): 0.5 (informal
rather than formalized). Principle 6 (local mediation through
cooperative): 1.0. D-type sum: 4.5/5.0.

\textbf{C principles:} Principle 3 (cooperative meetings with full
membership): 1.0. Principle 7 (Turkish government has not formally
recognized the local rules; legal status is ambiguous): 0.5. Principle 8
(no nested governance structure; single-level institution): 0.0. C-type
sum: 1.5/3.0.

\[D_{\text{norm}} = 90, \quad C_{\text{norm}} = 50, \quad B = \frac{2 \times 50}{90 + 50} \times 100 = 71.4\]

The B-score of 71.4 is above the \(1/\varphi\) threshold of 61.8,
consistent with the institution's documented sustainability. But the
C-deficit is visible: weak external recognition and no nested structure
make the institution vulnerable to disruption by external authority or
to scaling pressure. Ostrom notes exactly this: the Alanya system works
well at small scale but has not adapted successfully to regional
fisheries management.

\subsubsection{5.3 Open-Access Fishery (C-Dominant
Failure)}\label{open-access-fishery-c-dominant-failure}

Ostrom documents multiple pre-institutional or failed-institutional
cases that approximate open access: no boundary rules, no monitoring, no
sanctions. The Turkish inshore fisheries before Alanya developed its
system serve as a reference {[}Ost90{]}.

\textbf{D principles:} all effectively absent: no defined boundaries (1:
0.0), no locally matched rules (2: 0.0), no monitoring system (4: 0.0),
no sanctions (5: 0.0), no conflict resolution (6: 0.0). D-type sum:
0.0/5.0.

\textbf{C principles:} no formal collective-choice mechanism (3: 0.0),
no external recognition (7: 0.0), no nested structure (8: 0.0). C-type
sum: 0.0/3.0.

With both D and C at zero, the B-score formula is indeterminate (0/0).
In tholonic terms this represents a pre-institutional state: there is no
governance tholon yet. Individual extraction (ecological C) proceeds
without any governance D. The resource experiences C-dominant pressure
without any institutional N-formation. The trajectory is toward
depletion unless an institutional tholon self-organizes.

For comparison purposes: a minimal-D open access case (one partial
principle, no C) gives \(D_{\text{norm}} = 10\),
\(C_{\text{norm}} = 0\), \(B = 0\), still well below threshold.

\subsubsection{5.4 Atlantic Groundfish Regulatory Failure
(D-Dominant)}\label{atlantic-groundfish-regulatory-failure-d-dominant}

The Canadian and U.S. Atlantic groundfish (cod) collapse of 1992-1994 is
a documented case of centralized D-dominant governance failure. Federal
quota management had strong formal D but effectively excluded local user
participation and operated at a single national scale.

\textbf{D principles:} Principle 1 (licensing boundaries defined): 1.0.
Principle 2 (quotas set by federal scientists, partly matched to
conditions but miscalibrated due to data inadequacy): 0.5. Principle 4
(monitoring by federal observers): 1.0. Principle 5 (license suspension
and fines): 1.0. Principle 6 (federal administrative courts): 0.5.
D-type sum: 4.0/5.0.

\textbf{C principles:} Principle 3 (collective-choice: fisher advisory
committees existed but were consistently overridden; effectively 0.0
real influence): 0.0. Principle 7 (external recognition: the institution
was the state, so this principle is vacuously met but not meaningfully
applicable): 0.5. Principle 8 (nested enterprises: no nested
local-regional-national structure; monocentric federal management): 0.0.
C-type sum: 0.5/3.0.

\[D_{\text{norm}} = 80, \quad C_{\text{norm}} = 17, \quad B = \frac{2 \times 17}{80 + 17} \times 100 = 35.1\]

This B-score of 35.1, well below the 61.8 threshold, reflects a
D-dominant partial tholon. The institution had the structural properties
of rigidity: it maintained quotas that its own scientists were
questioning, overrode local knowledge about the declining state of the
stock, and collapsed when the mismatch between its rules and reality
became impossible to ignore. The collapse was sudden precisely because
D-dominant systems do not adapt gradually; they resist until structural
failure.

\begin{figure}
\centering
\pandocbounded{\includegraphics[keepaspectratio]{/home/jw/src/tv/docnav/Research/papers/17_ostrom-tholonic-governance/figures/17_bscore-cases.png}}
\caption{B-score comparison chart for the four case studies, showing
Törbel at B=100, Alanya at B=71.4, open-access at B=0, and Atlantic
groundfish at B=35, with the stability threshold at 61.8 marked.}
\end{figure}

\begin{center}\rule{0.5\linewidth}{0.5pt}\end{center}

\subsection{6. What the Tholonic Model
Adds}\label{what-the-tholonic-model-adds}

\subsubsection{6.1 From Binary Assessment to Continuous
Score}\label{from-binary-assessment-to-continuous-score}

Ostrom's framework assesses principle presence as binary: a principle is
either present (institutional success more likely) or absent (failure
more likely). Real governance institutions are almost never fully
described by this binary; principles are typically partially present,
partially enforced, or present in form but not in function.

The tholonic B-score provides a continuous metric. Alanya's B of 71.4 is
meaningfully different from Törbel's B of 100: it quantifies the gap and
identifies specifically which dimension (C-deficit from weak external
recognition and no nesting) accounts for it. This is not merely a
refinement in measurement precision; it enables differentiated
intervention targeting. An institution with \(B = 70\) from D-deficit
needs different interventions than one with \(B = 70\) from C-deficit,
even though both are above the threshold and both are technically
``sustainable.''

\subsubsection{6.2 Two Structurally Distinct Failure
Modes}\label{two-structurally-distinct-failure-modes}

Ostrom's framework has one primary failure mode: absence of sufficient
design principles. The tholonic model identifies two structurally
distinct failure modes with different profiles:

\textbf{C-dominant failure (tragedy of the commons):} \(D \approx 0\),
\(C > 0\). Resource is extracted without constraint. Collapse is
typically gradual and competitive: many small actors each extracting
more than their proportional share, with no mechanism for coordination.
The resource depletes steadily. Intervention requires D-construction:
boundary rules, monitoring, and sanctions.

\textbf{D-dominant failure (regulatory brittleness):} \(D > 0\),
\(C \approx 0\). Rules are enforced without participation. Collapse is
typically sudden and catastrophic: the institution holds its position
until reality overruns it. Intervention requires C-construction:
collective choice mechanisms, external legitimacy, and nested scales.
Paradoxically, this failure mode is harder to prevent because it looks
like success (compliance is maintained, rules are enforced) until near
the point of collapse.

The Atlantic groundfish case and Alanya's future vulnerability are both
D-dominant risks. Ostrom notes both but has no formal basis for
distinguishing them from C-dominant risks. The tholonic model does.

\subsubsection{6.3 Cross-Scale
Propagation}\label{cross-scale-propagation}

Because the tholonic model is recursive, a D-C imbalance at one level
propagates upward and downward through the governance hierarchy. An
institution that is D-dominant at the local level (Principle 3 absent:
no collective choice) undermines the N-state at the regional level
(Principle 8 absent: no nested integration), because the regional
institution cannot receive C-signals from the local level. Ostrom's
framework documents the nested enterprises principle (8) as important
but does not model the mechanism by which single-level failure
propagates across levels. The tholonic recursive hierarchy makes this
mechanism explicit: a partial tholon at level \(k\) deprives level
\(k+1\) of the child N it needs to form its own stable state.

\begin{center}\rule{0.5\linewidth}{0.5pt}\end{center}

\subsection{7. Falsifiable Predictions}\label{falsifiable-predictions}

The tholonic account of CPR governance generates three predictions that
are distinguishable from what Ostrom's framework predicts:

\textbf{Prediction 1: Two collapse signatures, not one.} C-dominant
governance failures (B near zero with D near zero) should show a
different temporal and distributional pattern than D-dominant governance
failures (B near zero with C near zero). C-dominant collapses should be
gradual, distributed across many small-scale actors, with continuous
resource decline. D-dominant collapses should be sudden, concentrated in
a short time window, and preceded by official reassurances of adequate
resource levels. This prediction can be tested against the historical
record of documented CPR collapses.

\textbf{Prediction 2: Continuous B-score predicts graded
sustainability.} Across Ostrom's documented case studies, institutions
with higher B-scores (regardless of whether the deficit is in D or C)
should show longer documented lifespans, lower frequency of crisis
episodes, and more rapid recovery from external disturbances.
Institutions with B near 61.8 (the boundary zone) should show unstable
oscillation: periods of function followed by crises, rather than either
stable operation or collapse. This prediction is testable against the
Ostrom case study archive.

\textbf{Prediction 3: D-C direction of deficit predicts intervention
effectiveness.} For institutions below the B-score threshold,
interventions that target the deficit side (C-deficit institutions
benefit from collective-choice investment; D-deficit institutions
benefit from monitoring and sanctions investment) should have higher
success rates than interventions that target the non-deficit side. This
is a specific prediction that goes against the common assumption that
all Ostrom principles are equally important in all contexts.

\begin{center}\rule{0.5\linewidth}{0.5pt}\end{center}

\subsection{8. Connection to the TVPCI Supply Chain
Framework}\label{connection-to-the-tvpci-supply-chain-framework}

The tholonic governance analysis developed here is structurally
continuous with the supply chain transparency work in papers
\href{https://github.com/baardev/truevalue/raw/main/docnav/Research/papers/2_supply-chain-transparency-tvpci/2_supply-chain-transparency-tvpci.pdf}{2}
and
\href{https://github.com/baardev/truevalue/raw/main/docnav/Research/papers/4_game-theoretic-triadic-balance/4_game-theoretic-triadic-balance.pdf}{4}
of this series
\href{https://github.com/baardev/truevalue/raw/main/docnav/Research/papers/2_supply-chain-transparency-tvpci/2_supply-chain-transparency-tvpci.pdf}{Mil26b},
\href{https://github.com/baardev/truevalue/raw/main/docnav/Research/papers/4_game-theoretic-triadic-balance/4_game-theoretic-triadic-balance.pdf}{Mil26d}.

A supply chain phase is a governance institution of a specific type: one
where the CPR is the production input (mineral stock, biomass, water),
the extractors are the supply chain actors at that phase, and the
governance institution is the combination of regulatory framework (D)
and industry cooperative mechanisms (C) that governs extraction and
processing at that phase.

The TVPCI transparency score
\href{https://github.com/baardev/truevalue/raw/main/docnav/Research/papers/2_supply-chain-transparency-tvpci/2_supply-chain-transparency-tvpci.pdf}{Mil26b}
measures the visibility of the custodial flow: how observable is the D-C
interaction at each phase? Phases with high TVPCI-D (well-documented
regulatory frameworks) and low TVPCI-C (poorly documented cooperative
participation) are precisely D-dominant partial tholons in the Ostrom
sense: they are brittle to regulatory change and resistant to adaptive
pressure from downstream actors. The Ostrom framework and the tholonic
supply chain analysis converge on the same diagnostic: governance health
requires balance between the constraining apparatus and the
participatory apparatus at every phase. The supply chain fails the same
way a CPR fails, for the same structural reason.

The connection also runs in the other direction: Ostrom's eight
principles, translated into supply chain language, provide a practical
checklist for diagnosing governance health at each supply chain phase. A
phase with clearly defined custody boundaries (Principle 1), processing
standards matched to local conditions (Principle 2), certified
monitoring (Principle 4), graduated compliance enforcement (Principle
5), dispute arbitration (Principle 6), producer representation in
standard-setting (Principle 3), regulatory recognition of cooperative
norms (Principle 7), and integration into multi-level certification
systems (Principle 8) scores \(B = 100\). Most commodity supply chain
phases score far below that, and the tholonic model specifies exactly
where the deficits lie and what their structural consequences are.

\begin{center}\rule{0.5\linewidth}{0.5pt}\end{center}

\subsection{9. Conclusions}\label{conclusions}

Elinor Ostrom's eight design principles for sustainable CPR governance
are not an arbitrary empirical list. They are the observational
fingerprint of a governance tholon: the structural conditions under
which a self-governing institution achieves and maintains the D-C
balance required to produce a stable N-state (sustainable resource
management). The five D-type principles constitute the constraining
apparatus of the institution; the three C-type principles constitute its
integrating and participatory apparatus. The tragedy of the commons is a
C-dominant partial tholon; D-dominant regulatory collapse is its
structural mirror image.

The tholonic model adds four things to Ostrom's framework: a continuous
balance score (\(B\)) that extends her binary principle-presence
assessment to a graded scalar; a formal account of two structurally
distinct failure modes with different temporal signatures; a mechanism
for cross-scale propagation through the recursive tholonic hierarchy;
and a direct connection to the supply chain governance analysis
developed elsewhere in this series.

Ostrom and the tholonic model are not competing frameworks. They are
complementary: Ostrom provides the empirical phenomenology of governance
health and failure; the tholonic model provides the structural account
of why those phenomena take the forms they do. Together they constitute
a more powerful diagnostic and predictive framework than either provides
alone.

\begin{center}\rule{0.5\linewidth}{0.5pt}\end{center}

\subsection{References}\label{references}

{[}Ost90{]} Ostrom, E. \emph{Governing the Commons: The Evolution of
Institutions for Collective Action.} Cambridge University Press, 1990.

{[}Ost05{]} Ostrom, E. \emph{Understanding Institutional Diversity.}
Princeton University Press, 2005.

{[}Ost10{]} Ostrom, E. ``Beyond Markets and States: Polycentric
Governance of Complex Economic Systems.'' \emph{American Economic
Review} 100:3 (2010), 641-672. (Nobel Prize lecture.)

{[}OstGar93{]} Ostrom, E., Gardner, R., and Walker, J. \emph{Rules,
Games, and Common-Pool Resources.} University of Michigan Press, 1993.

{[}Har68{]} Hardin, G. ``The Tragedy of the Commons.'' \emph{Science}
162:3859 (1968), 1243-1248.

{[}Dob14{]} Dietz, T., Ostrom, E., and Stern, P. C. ``The Struggle to
Govern the Commons.'' \emph{Science} 302:5652 (2003), 1907-1912.

\href{https://github.com/baardev/truevalue/raw/main/docnav/Research/papers/1_recursive-tholonic-five-constants/1_recursive-tholonic-five-constants.pdf}{Mil26a}
Milton, J. W. \emph{Recursive Tholonic Five Constants.} Clarity
Coalition,
\href{https://github.com/baardev/truevalue/raw/main/docnav/Research/papers/1_recursive-tholonic-five-constants/1_recursive-tholonic-five-constants.pdf}{paper
1}, 2026.

\href{https://github.com/baardev/truevalue/raw/main/docnav/Research/papers/2_supply-chain-transparency-tvpci/2_supply-chain-transparency-tvpci.pdf}{Mil26b}
Milton, J. W. \emph{Supply Chain Transparency and the TVPCI.} Clarity
Coalition,
\href{https://github.com/baardev/truevalue/raw/main/docnav/Research/papers/2_supply-chain-transparency-tvpci/2_supply-chain-transparency-tvpci.pdf}{paper
2}, 2026.

\href{https://github.com/baardev/truevalue/raw/main/docnav/Research/papers/3_minimal-recursive-triadic-framework/3_minimal-recursive-triadic-framework.pdf}{Mil26c}
Milton, J. W. \emph{Minimal Recursive Triadic Framework.} Clarity
Coalition,
\href{https://github.com/baardev/truevalue/raw/main/docnav/Research/papers/3_minimal-recursive-triadic-framework/3_minimal-recursive-triadic-framework.pdf}{paper
3}, 2026.

\href{https://github.com/baardev/truevalue/raw/main/docnav/Research/papers/4_game-theoretic-triadic-balance/4_game-theoretic-triadic-balance.pdf}{Mil26d}
Milton, J. W. \emph{Game-Theoretic Triadic Balance.} Clarity Coalition,
\href{https://github.com/baardev/truevalue/raw/main/docnav/Research/papers/4_game-theoretic-triadic-balance/4_game-theoretic-triadic-balance.pdf}{paper
4}, 2026.

\href{https://github.com/baardev/truevalue/raw/main/docnav/Research/papers/16_tholonic-spinoza-leibniz/16_tholonic-spinoza-leibniz.pdf}{Mil26e}
Milton, J. W. \emph{Tholonic N-D-C Structure and Its Philosophical
Predecessors: Spinoza and Leibniz.} Clarity Coalition,
\href{https://github.com/baardev/truevalue/raw/main/docnav/Research/papers/16_tholonic-spinoza-leibniz/16_tholonic-spinoza-leibniz.pdf}{paper
16}, 2026.

{[}Cox10{]} Cox, M., Arnold, G., and Villamayor Tomás, S. ``A Review of
Design Principles for Community-Based Natural Resource Management.''
\emph{Ecology and Society} 15:4 (2010), article 38.

{[}Rod20{]} Rodríguez-Meza, J. \emph{The Eight Ostrom Principles in
Practice: Cross-Country Evidence.} Working paper, 2020.

\end{document}
